\scanpage{39}
\begin{center}
UNIVERSITÉ DE PARIS

\figure{\begin{tikzpicture}\draw (-.4,0)--(-.06,0);\fill (0,0) circle (.7pt);\draw (.06,0)--(.4,0);\end{tikzpicture}}

ÉCOLE NORMALE SUPÉRIEURE

\figure{\begin{tikzpicture}\draw (-.4,0)--(-.06,0);\fill (0,0) circle (.7pt);\draw (.06,0)--(.4,0);\end{tikzpicture}}

45, RUE D'ULM

ODÉON 06$\cdot$45

\figure{\begin{tikzpicture}\draw (-.4,0)--(-.06,0);\fill (0,0) circle (.7pt);\draw (.06,0)--(.4,0);\end{tikzpicture}}
\end{center}

\deleted{\unsure{En particulier si $E$ est}}

\uuline{Th 2} : Si $G$ est un groupe fini d'ordre $n$, l'image de
\[
p_*:H_*(E)\longrightarrow H_*(B)
\]
contient tous les éléments \deleted{\unsure{d'ordre $n$}} divisibles par $n$ dans $H_*(B)$.

\uuline{Dém} : Soit $\zeta$ \deleted{\unsure{un cycle}}$=\sum_{i\in I}a_i\sigma_i$ un cycle de $B$, où $\sigma_i$ sont des simplexes.
Pour tout $\sigma_i$, soit $\tau_i$ simplexe de $E$ tel que $p\tau_i=\sigma_i$. Alors
\[
\theta=\sum_{g\in G,\,i\in I}ga_i\tau_i
\]
est un cycle de $E$, et $p(\theta)=\deleted{\unsure{\zeta}}n\zeta$ \deleted{\illegible{1}}. D'où le Th 2.

\uuline{Th 3} : Si \deleted{\unsure{l'espace}} $E$ est \deleted{\illegible{1}} contractile, $S=$ ss-groupe engendré par les stabilisateurs fini d'ordre $n$, $\pi_i(B)$ est un groupe \deleted{\unsure{fini}} dont \inserted{tous les éléments sont} d'ordre qui divise une puissance de $n$ pour $i>1$

\uuline{Dém} : \deleted{\unsure{D'après}} D'après le Th. 1, $E/S$ est le revêtement universel de $B$, d'où $\pi_i(B)=\pi_i(E/S)$. D'après le Th. 2,

\deleted{\unsure{tous}} les $H_i(E/S)$

\mspage{38}

